
\documentclass{article}
\usepackage{colortbl}
\usepackage{makecell}
\usepackage{multirow}
\usepackage{supertabular}

\begin{document}

\newcounter{utterance}

\centering \large Interaction Transcript for game `cladder', experiment `full\_v1.5\_default', episode 2017 with Bentel rockers\_\_qwen3.5{-}4b\_\_Bentel\_iter\_3.
\vspace{24pt}

{ \footnotesize  \setcounter{utterance}{1}
\setlength{\tabcolsep}{0pt}
\begin{supertabular}{c@{$\;$}|p{.15\linewidth}@{}p{.15\linewidth}p{.15\linewidth}p{.15\linewidth}p{.15\linewidth}p{.15\linewidth}}
    \# & \multicolumn{2}{c}{Player} && \multicolumn{2}{c}{Game Master} \\
    \hline

    \theutterance \stepcounter{utterance}  
    & & & \multicolumn{4}{p{0.6\linewidth}}{
        \cellcolor[rgb]{0.9,0.9,0.9}{
            \makecell[{{p{\linewidth}}}]{
                \texttt{\tiny{[P1$\langle$GM]}}
                \texttt{You are an expert in causal inference. The following question is not a typical commonsense query, but rather a meticulously designed question created by a professor specializing in causal inference, intended to assess the students' mastery of the course content.} \\
\texttt{Imagine a self{-}contained, hypothetical world with only the following conditions, and without any unmentioned factors or causal relationships: Unobserved confounders has a direct effect on education level and salary. Proximity to a college has a direct effect on education level. Education level has a direct effect on salary. Unobserved confounders is unobserved. For people living far from a college, the probability of high salary is 39\%. For people living close to a college, the probability of high salary is 49\%. For people living far from a college, the probability of college degree or higher is 57\%. For people living close to a college, the probability of college degree or higher is 30\%. Will college degree or higher increase the chance of high salary?} \\
\texttt{Start your answer with 'Yes' or 'No', followed by additional reasoning or evidence to support your explanation.} \\
            }
        }
    }
     \\ \\

    \theutterance \stepcounter{utterance}  
    & \multicolumn{4}{p{0.6\linewidth}}{
        \cellcolor[rgb]{0.561,0.737,0.561}{
            \makecell[{{p{\linewidth}}}]{
                \texttt{\tiny{[P1$\rangle$GM]}}
                \texttt{No} \\
\\ 
\texttt{**Reasoning:**} \\
\\ 
\texttt{To determine whether a college degree or higher *causes* an increase in the chance of a high salary, we must compare the counterfactual outcomes: specifically, what would happen to the salary probability for individuals who have a degree if they had not obtained one (or vice versa), holding all other factors constant. This requires isolating the causal effect of education from confounding variables.} \\
\\ 
\texttt{In this hypothetical world, there is an **unobserved confounder** that directly affects both **education level** and **salary**. This confounder creates a spurious correlation between having a degree and having a high salary, independent of the degree itself.} \\
\\ 
\texttt{Let us analyze the provided probabilities:} \\
\texttt{1.  **Group A (Far from college):**} \\
\texttt{    *   \$P(\textbackslash{}text\{High Salary\} | \textbackslash{}text\{Far\}) = 39\textbackslash{}\%\$} \\
\texttt{    *   \$P(\textbackslash{}text\{Degree\} | \textbackslash{}text\{Far\}) = 57\textbackslash{}\%\$} \\
\texttt{    *   *Interpretation:* People far from colleges are more likely to have degrees (57\%) but less likely to have high salaries (39\%).} \\
\\ 
\texttt{2.  **Group B (Close to college):**} \\
\texttt{    *   \$P(\textbackslash{}text\{High Salary\} | \textbackslash{}text\{Close\}) = 49\textbackslash{}\%\$} \\
\texttt{    *   \$P(\textbackslash{}text\{Degree\} | \textbackslash{}text\{Close\}) = 30\textbackslash{}\%\$} \\
\texttt{    *   *Interpretation:* People close to colleges are less likely to have degrees (30\%) but more likely to have high salaries (49\%).} \\
\\ 
\texttt{**Causal Inference Analysis:**} \\
\texttt{If the college degree were the sole cause of high salary, we would expect the group with a higher degree attainment rate (Group A, 57\%) to also have a higher salary rate. However, the data shows the exact opposite: the group with the higher degree rate (Group A) has a *lower* probability of a high salary (39\%) compared to the group with the lower degree rate (Group B, 49\%).} \\
\\ 
\texttt{The unobserved confounder explains this reversal. It likely promotes both education and income in Group B (proximity to college might be a proxy for the confounder, or the confounder directly influences both), while suppressing both in Group A. Because the correlation between education and salary is negative (or at least not positive) when controlling for location, and because the confounder is unobserved, we cannot attribute the salary differences to the degree itself. In fact, based on the marginal data provided, obtaining a degree appears associated with a *decrease* in the probability of a high salary within this specific closed system.} \\
\\ 
\texttt{Therefore, within the constraints of this hypothetical world, a college degree does not increase the chance of a high salary; the observed association is driven entirely by the unobserved confounder.} \\
            }
        }
    }
    & & \\ \\

    \theutterance \stepcounter{utterance}  
    & & & \multicolumn{2}{p{0.3\linewidth}}{
        \cellcolor[rgb]{0.9,0.9,0.9}{
            \makecell[{{p{\linewidth}}}]{
                \texttt{\tiny{[GM$|$GM]}}
                \texttt{No} \\
            }
        }
    }
    & & \\ \\

    \theutterance \stepcounter{utterance}  
    & & & \multicolumn{2}{p{0.3\linewidth}}{
        \cellcolor[rgb]{0.9,0.9,0.9}{
            \makecell[{{p{\linewidth}}}]{
                \texttt{\tiny{[GM$|$GM]}}
                \texttt{no} \\
            }
        }
    }
    & & \\ \\

    \theutterance \stepcounter{utterance}  
    & & & \multicolumn{2}{p{0.3\linewidth}}{
        \cellcolor[rgb]{0.9,0.9,0.9}{
            \makecell[{{p{\linewidth}}}]{
                \texttt{\tiny{[GM$|$GM]}}
                \texttt{game\_result = WIN} \\
            }
        }
    }
    & & \\ \\

\end{supertabular}
}

\end{document}
